\chapter{Especificação de Requisitos}\label{cap:especificacao}
% ===========================================================================

Este capítulo apresenta os requisitos da solução proposta para uma livraria
digital em modelo de catálogo, que utiliza um algoritmo \emph{oblivious} para
viabilizar a venda dos itens mantendo a privacidade sobre qual item foi
comprado.

\section{Requisitos Funcionais}

\begin{itemize}
\item \textbf{RF01 --- Cadastro de administradores:} o sistema deve permitir o
cadastro, autenticação e gerenciamento de perfis de administradores.

\item \textbf{RF02 --- Manuseio de livros no catálogo:} o sistema deve
permitir ao administrador cadastrar, editar, consultar e remover livros do
catálogo.

\item \textbf{RF03 --- Política de preços discretos:} o sistema deve permitir
a precificação de cada livro dentro de um conjunto de preços pré-definidos.

\item \textbf{RF04 --- Busca e navegação no catálogo:} o sistema deve permitir
que clientes pesquisem livros por critérios como título, autor, categoria e
preço. O catálogo é publicado pelo sistema web como um pacote completo, na
mesma forma para todos os clientes, e a busca e a navegação são feitas no
programa local, de modo que o servidor não observe os critérios de busca nem
os livros consultados. O sistema web não oferece vitrine pública nem imagens
de capa avulsas; sua única página pública é a de apresentação e download, igual
para todos os visitantes.

\item \textbf{RF05 --- Compra sem cadastro de cliente:} o sistema deve
permitir compra sem criação de conta de cliente.

\item \textbf{RF06 --- E-mail obrigatório na compra:} no fluxo de compra, o
cliente deve informar um e-mail válido para recebimento da resposta do
protocolo de OT, a partir da qual o programa local extrai a chave do produto
digital.

\item \textbf{RF07 --- Compra privada de N itens em grupo:} o cliente deve
conseguir comprar N livros dentre um grupo maior de livros com o mesmo preço
sem que a informação dos itens específicos comprados seja enviada ao servidor
em formato que ele consiga decifrar. Cada grupo possui uma versão imutável
(composição e ordem dos livros), e a compra referencia essa versão.

\item \textbf{RF08 --- Entrega de material criptográfico para cliente:} o
sistema deve fornecer ao cliente os dados necessários (a resposta do
protocolo de OT e informações relacionadas) para que as chaves simétricas dos
livros escolhidos sejam extraídas em ambiente local. O servidor não conhece as
chaves finais obtidas pelo cliente.

\item \textbf{RF09 --- Arquitetura com três componentes:} a solução deve ser
composta por (i) um sistema web para administração e publicação do catálogo,
pagamento, execução do OT e assinatura cega de tokens; (ii) uma CDN, que
armazena os livros cifrados e os entrega mediante token válido; e (iii) um
programa local, instalado no computador do cliente e não fornecido pelo
sistema web em tempo de execução, que permite navegar no catálogo, guarda o
estado secreto do cliente, comunica-se apenas com o sistema web e com a CDN e
descriptografa os livros.

\item \textbf{RF10 --- Descriptografia local:} o cliente deve conseguir
descriptografar o item resgatado por meio do programa local, a partir do livro
cifrado e da resposta do OT. A chave do livro é obtida e usada apenas no
computador do cliente e nunca é transmitida pela rede. O programa local
comunica-se exclusivamente com o sistema web e com a CDN, sem telemetria nem
serviços de terceiros.

\item \textbf{RF11 --- Gestão do segredo no programa local:} o programa local
deve gerar a requisição de OT a partir de uma semente secreta, obtida de um
gerador criptograficamente seguro do sistema, com pelo menos 128 bits, nova a
cada compra e nunca reutilizada. O programa guarda a semente e as escolhas,
necessárias para extrair, a partir da resposta do servidor, as chaves
criptográficas dos itens adquiridos, inclusive a partir da resposta recebida
por e-mail (RF14) caso a compra seja interrompida após o pagamento. A perda da
semente implica a perda de acesso aos itens daquela compra.

\item \textbf{RF12 --- Não acesso a itens não comprados:} o cliente não deve
conseguir descriptografar livros que não tenham sido comprados.

\item \textbf{RF13 --- Processamento de pagamento:} o sistema deve realizar o
processamento do pagamento da compra por meio de mecanismo integrado de
cobrança (gateway de pagamento) e retornar o status da transação ao cliente.

\item \textbf{RF14 --- Envio da resposta do OT por e-mail:} após a confirmação
do pagamento, o sistema deve enviar para o e-mail informado a resposta do OT
correspondente à compra. O e-mail não contém a chave do livro, que o servidor
desconhece.

\item \textbf{RF15 --- Persistência para reenvio ao cliente:} o sistema deve
manter os dados da compra de forma que o cliente possa solicitar o reenvio da
resposta do OT para o mesmo e-mail informado originalmente, dentro de um prazo
definido pela plataforma.

\item \textbf{RF16 --- Emissão de tokens com assinatura cega:} para cada
compra de N itens, o sistema web deve assinar às cegas N valores gerados pelo
cliente, resultando em N tokens que autorizam, cada um, o download de um livro
na CDN. Os tokens não devem poder ser ligados à sessão de assinatura.

\item \textbf{RF17 --- Entrega de livros pela CDN mediante token:} a CDN deve
entregar um livro cifrado apenas mediante a apresentação de um token com
assinatura válida e ainda não utilizado.

\item \textbf{RF18 --- Reporte de uso da CDN:} a CDN deve reportar ao sistema
web os tokens utilizados e a quantidade de livros entregues, para fins de
remuneração. O reporte não deve conter o livro associado a cada token nem
horários que permitam correlação com as compras.
\end{itemize}

\section{Requisitos Não Funcionais}

\begin{itemize}
\item \textbf{RNF01 --- Segurança criptográfica:} o mecanismo de proteção deve
atingir um nível mínimo de 128 bits de segurança, isto é, um ataque ao
mecanismo deve exigir da ordem de $2^{128}$ operações.

\item \textbf{RNF02 --- Tempo limite de busca:} cada operação de busca no
catálogo deve ser concluída em um tempo de resposta limite de 5 segundos.

\item \textbf{RNF03 --- Tempo limite de finalização da compra:} a finalização
da compra deve ser concluída em um tempo de resposta limite de 1 minuto.

\item \textbf{RNF04 --- Tempo limite de envio por e-mail:} o sistema deve
enviar o e-mail contendo a resposta do OT dentro de um tempo de resposta
limite de 1 minuto após confirmação do pagamento.

\item \textbf{RNF05 --- Capacidade de processamento:} a plataforma deve
suportar uma quantidade mínima de 600 transações por minuto.

\item \textbf{RNF06 --- Disponibilidade:} a plataforma deve atingir um valor
mínimo de 99\% de \emph{uptime} por mês.

\item \textbf{RNF07 --- Usabilidade:} o fluxo de compra deve atingir uma taxa
mínima de conclusão sem suporte externo igual a 99\%.

\item \textbf{RNF08 --- Manutenibilidade:} alterações evolutivas de baixa
complexidade devem ser implementadas em um tempo máximo de 8 horas úteis de
manutenção.
\end{itemize}

\section{Modelo de Ameaça}

O projeto adota as seguintes premissas de segurança, que orientam as decisões
de arquitetura:

\begin{itemize}
\item \textbf{Site e CDN não são coniventes.} Essa premissa é garantida por
auditoria. Com isso, o site sabe \emph{quem} comprou, mas não \emph{o quê}; a
CDN sabe \emph{o quê} foi baixado, mas não \emph{quem} comprou. Essa separação
de conhecimento é intencional e não constitui uma falha do sistema.

\item \textbf{O reporte da CDN ao site deve ser mínimo.} Ele contém apenas os
tokens utilizados e a quantidade de entregas, nunca o livro associado a cada
token nem horários que permitam correlacionar downloads com compras (RF18).

\item \textbf{Os tokens não podem ser ligados à compra.} A assinatura cega
garante que o site, ao receber um token no reporte da CDN, não consiga
associá-lo à sessão de assinatura em que foi emitido (RF16).

\item \textbf{O servidor não conhece a chave final do cliente.} O site executa
sua parte do OT sem saber quais chaves o cliente extraiu. Por isso o e-mail
(RF14) e o reenvio (RF15) carregam a \emph{resposta do OT}, e não a chave.
Somente o programa local, que guarda o segredo da requisição (RF11), consegue
extrair a chave.

\item \textbf{O código do receptor do OT não vem do servidor.} A requisição de
OT é gerada pelo programa local, instalado pelo cliente, e não por código
entregue pelo sistema web durante a compra. Assim, a privacidade da escolha
(RF07) decorre do protocolo de OT e não depende de o servidor fornecer código
honesto.

\item \textbf{Navegação e busca não revelam interesse.} O programa local
baixa o catálogo inteiro, com miniaturas de capa de tamanho fixo, e faz a
busca e a navegação localmente (RF04); o sistema web observa apenas que o
pacote completo foi baixado. Como o código de navegação não é fornecido pelo
sistema web, a proteção não depende de o servidor entregar código honesto.
Trata-se de uma medida de engenharia, não de uma garantia do protocolo;
complementam-na a ausência de vitrine pública, de imagens avulsas e de
registros de acesso no sistema web.

\item \textbf{O programa local fala apenas com o sistema web e com a CDN.} A
chave final e a semente nunca deixam o computador do cliente (RF10, RF11). O
endereço de rede do cliente é visível à CDN; a ligação entre esse dado e a
identidade do comprador exigiria conivência entre site e CDN, excluída pela
premissa acima.

\item \textbf{O cegamento dos tokens é removido no programa local.} Os valores
cegados são gerados e descegados no programa local, e não no navegador, pelo
mesmo motivo do item sobre o código do receptor.
\end{itemize}

\section{Casos de uso e diagrama de fluxo}

\begin{table}[htbp]
\centering
\caption{Agentes e Casos de Uso do Sistema}\label{tab:casos-de-uso}
\small
\begin{tabular}{@{}l l p{0.5\textwidth}@{}}
\toprule
\textbf{ID} & \textbf{Agente} & \textbf{Caso de Uso} \\
\midrule
UC01 & Administrador & Gerenciar conta de administrador (cadastro, login e edição) \\
UC02 & Administrador & Gerenciar catálogo de livros (cadastrar, editar, consultar e remover) \\
UC03 & Administrador & Definir política de preços discretos para os livros \\
UC04 & Cliente (Programa Local) & Buscar e navegar no catálogo de livros \\
UC05 & Cliente & Realizar compra sem necessidade de cadastro \\
UC06 & Cliente & Informar e-mail válido para recebimento da resposta do OT \\
UC07 & Cliente & Realizar compra privada de N itens utilizando protocolo \emph{oblivious transfer} \\
UC08 & Sistema & Processar pagamento da compra \\
UC09 & Sistema & Executar o OT e entregar a resposta ao cliente \\
UC10 & Sistema & Enviar resposta do OT por e-mail após pagamento \\
UC11 & Cliente & Solicitar reenvio da resposta do OT \\
UC12 & Sistema & Persistir dados da compra para reenvio \\
UC13 & Sistema & Assinar às cegas os N valores do cliente (emissão de tokens) \\
UC14 & Cliente & Baixar livro cifrado na CDN apresentando um token (via programa local) \\
UC15 & CDN & Validar token e entregar livro cifrado \\
UC16 & CDN & Reportar tokens utilizados e quantidade ao site \\
UC17 & Sistema & Validar tokens reportados e remunerar a CDN \\
UC18 & Cliente (Programa Local) & Descriptografar conteúdo localmente \\
UC19 & Sistema (Programa Local) & Guardar o segredo da requisição e extrair a chave localmente \\
UC20 & Sistema & Garantir que itens não comprados não possam ser acessados \\
\bottomrule
\end{tabular}
\fonte{Elaborado pelos autores.}
\end{table}

A \autoref{fig:seq-compra} apresenta o fluxo da administração do catálogo até
a compra, e a \autoref{fig:seq-entrega} apresenta a entrega pela CDN, o
reporte de uso e a descriptografia local, incluindo a retomada de uma compra
interrompida após o pagamento.

% ---------------------------------------------------------------------------
\begin{figure}[htbp]
\centering
\caption{Diagrama de sequência: administração, navegação e compra}\label{fig:seq-compra}
\begin{tikzpicture}
  \sdlifelines{-14.6}
  \sdheads

  \sdbanner{-1.0}{Administração}
  \sdcall{\xAdm}{\xWeb}{-1.6}{Gerenciar catálogo}
  \sdcall{\xAdm}{\xWeb}{-2.2}{Definir preços discretos}

  \sdbanner{-2.9}{Navegação}
  \sdcall{\xCli}{\xPrg}{-3.6}{Buscar livros}
  \sdcall{\xPrg}{\xWeb}{-4.15}{Baixar pacote do catálogo}
  \sdreply{\xWeb}{\xPrg}{-4.7}{Pacote completo (igual para todos)}
  \sdselfL{\xPrg}{-4.95}{Filtrar localmente}

  \sdbanner{-5.75}{Compra}
  \sdcall{\xCli}{\xPrg}{-6.35}{Escolher grupo (versão) e N livros do grupo}
  \sdselfL{\xPrg}{-6.75}{Gerar semente, requisição OT e N valores\\cegados; guardar semente e escolhas}
  \sdcall{\xPrg}{\xWeb}{-7.8}{E-mail + requisição OT + N valores cegados}
  \sdcall{\xWeb}{\xGat}{-8.35}{Processar pagamento (grupo $\times$ N)}
  \sdreply{\xGat}{\xWeb}{-8.9}{Status da transação}

  \sdframe{-0.8}{-9.25}{13.3}{-14.25}{alt}{Pagamento aprovado}
  \sdselfR{\xWeb}{-9.85}{Calcular resposta do OT (sem saber a escolha)}
  \sdselfR{\xWeb}{-10.45}{Assinar às cegas os N valores}
  \sdreply{\xWeb}{\xPrg}{-11.25}{Resposta do OT + N assinaturas cegas}
  \sdcall{\xWeb}{\xCli}{-11.75}{Enviar resposta do OT (e-mail)}
  \sdselfR{\xWeb}{-12.15}{Persistir resposta do OT para reenvio}
  \sdselfL{\xPrg}{-12.75}{Remover cegamento: N tokens}
  \sdelse{-0.8}{13.3}{-13.3}{Pagamento recusado}
  \sdreply{\xWeb}{\xPrg}{-13.9}{Informar falha}
\end{tikzpicture}
\fonte{Elaborado pelos autores.}
\end{figure}

% ---------------------------------------------------------------------------
\begin{figure}[htbp]
\centering
\caption{Diagrama de sequência: entrega, reporte da CDN e descriptografia local}\label{fig:seq-entrega}
\begin{tikzpicture}
  \sdlifelines{-15.4}
  \sdheads

  \sdbanner{-1.0}{Entrega}
  \sdframe{6.5}{-1.4}{13.3}{-3.8}{loop}{Para cada livro desejado (até N)}
  \sdcall{\xPrg}{\xCdn}{-2.3}{Token + livro solicitado}
  \sdselfL{\xCdn}{-2.6}{Validar assinatura e\\uso único do token}
  \sdreply{\xCdn}{\xPrg}{-3.4}{Livro cifrado}

  \sdbanner{-4.5}{Reporte da CDN}
  \sdcall{\xCdn}{\xWeb}{-5.1}{Tokens utilizados + quantidade (sem livro nem horário)}
  \sdselfR{\xWeb}{-5.5}{Validar tokens}
  \sdframe{3.6}{-6.1}{11.0}{-8.55}{alt}{Tokens válidos}
  \sdcall{\xWeb}{\xCdn}{-6.85}{Efetuar pagamento pelos downloads}
  \sdelse{3.6}{11.0}{-7.25}{Tokens inválidos}
  \sdcall{\xWeb}{\xCdn}{-8.2}{Rejeitar pagamento / sinalizar erro}

  \sdbanner{-9.0}{Descriptografia (local)}
  \sdframe{-0.8}{-9.4}{13.3}{-11.1}{opt}{Compra interrompida após o pagamento}
  \sdcall{\xCli}{\xPrg}{-10.15}{Resposta do OT recebida por e-mail}
  \sdselfL{\xPrg}{-10.5}{Reproduzir o receptor a partir da semente}
  \sdselfL{\xPrg}{-11.6}{Extrair as chaves com a semente e a resposta do OT}
  \sdselfL{\xPrg}{-12.2}{Descriptografar livro}
  \sdframe{-0.8}{-12.8}{13.3}{-15.0}{alt}{Livro comprado}
  \sdreply{\xPrg}{\xCli}{-13.55}{Exibir conteúdo}
  \sdelse{-0.8}{13.3}{-13.95}{Livro não comprado}
  \sdreply{\xPrg}{\xCli}{-14.7}{Acesso negado}
\end{tikzpicture}
\fonte{Elaborado pelos autores.}
\end{figure}
